\chapter{The Metric Under the Microscope}
\label{ch:metric}

Competitions are optimization problems, and the objective is the metric ---
not accuracy in the abstract, not your intuition of quality, but one specific
formula evaluated on one specific table. Small print in that formula changes
what models you should build, what loss you should train with, and even what
post-processing is worth doing. This chapter dissects the Jane Street metric
in detail and extracts the general lessons.

\section{The zero-mean weighted \texorpdfstring{$R^2$}{R2}}

The competition scores predictions $\hat y_i$ against realized values $y_i$
with per-row weights $w_i$:
\begin{equation}
\Rw \;=\; 1 - \frac{\sum_i w_i\,(y_i - \hat y_i)^2}{\sum_i w_i\, y_i^2}.
\label{eq:wr2}
\end{equation}

Three deliberate details distinguish this from the textbook $R^2$, and each
one carries a lesson.

\subsection{The denominator is \emph{not} centered}

The usual coefficient of determination compares your errors to those of the
sample-mean predictor: denominator $\sum (y_i - \bar y)^2$. Here the
denominator is $\sum w_i y_i^2$ --- the comparison point is the
\textbf{all-zeros predictor}. Consequences:

\begin{itemize}
\item \textbf{Zero is the origin of the score.} Predicting $\hat y \equiv 0$
  scores exactly $\Rw = 0$. Any model must beat ``do nothing.'' In trading
  terms: the benchmark is staying flat, which is the honest benchmark for a
  signal that drives position-taking.
\item \textbf{Predicting the historical mean can score \emph{negative}.}
  If the target's mean in your training window differs from zero (drift!),
  carrying that mean forward adds $(\bar y_{\text{train}})^2$ to your error
  with no compensation in the denominator. At our SNR, even tiny biases
  matter: a constant offset $b$ costs approximately $b^2 / \overline{y^2}$
  of score. With $\overline{y^2} \approx 0.6$ and target scores of $0.01$,
  a bias of $b = 0.05$ burns $0.004$ --- nearly half a good model's entire
  score. \emph{Check your prediction mean.}
\item \textbf{The metric rewards shrinkage.} Decompose the score for a
  prediction $\hat y = \alpha f$, where $f$ is your (imperfect) signal
  estimate and $\alpha$ a scale you control. The optimal $\alpha$ is the
  regression coefficient of $y$ on $f$:
  $\alpha^\star = \Cov_w(y, f)/\Var_w(f)$, and \emph{underconfident
  predictions lose less than overconfident ones}. When a model chronically
  overpredicts magnitudes (common after training on MSE with outliers),
  a global shrink toward zero, fit on validation, is free money. Our
  pipeline clips final predictions to $[-5, 5]$ for the same tail-risk
  reason: one insane prediction on a high-weight row can erase a day of
  signal.
\end{itemize}

\subsection{The weights are part of the objective}

Every row carries a weight $w_i$ (in this competition, plausibly
proportional to a liquidity or dollar-impact measure). The metric is a
weighted sum, which means \textbf{training must be weighted identically}:
the loss on a training example must be multiplied by the same $w_i$ used at
evaluation. This sounds obvious; in practice it is one of the top-five most
common silent errors in public notebooks. Symptoms include: a model that
looks fine on unweighted MSE but underperforms on the leaderboard, because
it spends capacity on low-weight rows.

Two secondary effects of weighting deserve attention:

\begin{itemize}
\item \textbf{Effective sample size shrinks.} With weights of varying
  magnitude, the variance-reducing power of $n$ rows is
  $n_{\text{eff}} = (\sum w_i)^2 / \sum w_i^2 \le n$. Heavy-tailed weights
  concentrate the metric on fewer effective rows, making validation noisier
  than row counts suggest --- one more reason for the replication discipline
  of Chapter~\ref{ch:problem}.
\item \textbf{Weighted early stopping.} If your training framework's
  early-stopping metric ignores the weights (XGBoost's default RMSE on an
  unweighted eval set, for instance), you are selecting the stopping point
  for the wrong objective. Pass the weights to the eval set too
  (\code{sample\_weight\_eval\_set} in the sklearn wrapper).
\end{itemize}

\subsection{It is a squared-error metric --- with everything that implies}

$\Rw$ is affine in weighted MSE, so the Bayes-optimal prediction is the
conditional mean $\E[y \mid \text{features}]$, not the median, not a
quantile, and not a direction. Squared error punishes large misses
quadratically, which at heavy-tailed targets means \textbf{tail rows
dominate gradients}. Practical mitigations, all used in our campaign or by
the solutions we studied:

\begin{itemize}
\item \textbf{Clip the target's influence}, either by clipping predictions
  (we clip to $[-5,5]$) or by gradient clipping during training
  (global-norm 1.0 on all neural models).
\item \textbf{Robustify the inputs}, not the loss: quantile-clip features at
  fitted $[0.1\%, 99.9\%]$ bounds (Chapter~\ref{ch:preprocessing}) so
  extreme feature rows cannot generate extreme predictions.
\item \textbf{Consider a correlation term in the training loss.} The DRW
  crypto winner trained his MLP on
  $0.6\cdot\mathrm{MSE} + 0.4\cdot(1-\rho_{\mathrm{Pearson}})$ even though
  the competition metric was correlation-adjacent. His stated reason
  generalizes: at low SNR, pure MSE gradients are dominated by irreducible
  noise, and a correlation term keeps optimization pointed at
  \emph{co-movement} --- the part you can actually influence. When the
  competition metric is exactly $\Rw$, training directly on
  $1-\Rw$ per batch (as our \code{WeightedR2Loss} does) is the cleanest
  choice; the denominator per batch also acts as a natural per-day
  normalization when each batch is one trading day.
\end{itemize}

\section{Training loss versus evaluation metric}

A recurring genre decision: do you train on the metric itself, or on a proxy?
Our answer, empirical and consistent with both reference solutions:

\begin{center}
\begin{tabularx}{\textwidth}{@{}lXX@{}}
\toprule
model family & train loss & rationale \\
\midrule
gradient-boosted trees & plain weighted MSE & tree boosting needs pointwise
gradients/hessians; the $\Rw$ affine shift does not change split decisions;
weighting enters via sample weights. \\
neural sequence models & $1-\Rw$ per day-batch & each batch is one full
trading day, so the batch denominator $\sum w y^2$ is a stable day-level
normalizer; the loss is then scale-calibrated across days with different
volatility. \\
MLP (DRW winner) & $0.6\,\mathrm{MSE} + 0.4\,(1-\rho)$ & correlation term
stabilizes low-SNR optimization; MSE term keeps calibration. \\
\bottomrule
\end{tabularx}
\end{center}

\begin{keybox}[: batch = one day makes the metric trainable]
The per-day batch convention (Chapter~\ref{ch:rnn}) has a quiet side effect:
it makes $1-\Rw$ a well-behaved training loss. Computed over one day's rows,
the denominator is a meaningful quantity (that day's weighted target energy),
and each day contributes one loss value of comparable scale regardless of
that day's volatility. Random-row batches would make the batch-$\Rw$
denominator noisy and the loss erratic. Metric-aware batching is an
underrated trick.
\end{keybox}

\section{What ``good'' looks like on this metric: a calibration table}

Because almost nobody has intuitions for scores of $0.01$, it is worth
building a mental ladder. All numbers are weighted $R^2$ on our 100-day
validation tail unless stated:

\begin{center}
\begin{tabular}{@{}lr@{}}
\toprule
predictor & $\Rw$ \\
\midrule
all zeros & $0.0000$ \\
carry forward training mean & $\approx -0.001$ \\
plain MLP on all features & $-0.0026$ \\
10-feature linear probe & $+0.002$ to $+0.004$ \\
well-tuned XGBoost, 200 recent days & $+0.0059$ \\
XGBoost + engineered/symbolic features & $+0.0064$ \\
GRU with aux heads, static & $+0.0076$ \\
same, with daily online refit & $+0.0089$ \\
LSTM with aux heads, online, tuned refit LR & $+0.0092$ \\
3-model decorrelated ensemble & $+0.0104$ \\
8th place (private LB, different data) & $+0.0104$ \\
1st place (private LB) & $+0.0139$ \\
\bottomrule
\end{tabular}
\end{center}

Two readings of this ladder are worth internalizing. First, the \emph{entire
competitive range} --- from ``did something sensible'' to ``best on Earth''
--- spans about $0.012$. Steps you might consider negligible elsewhere
($+0.0005$) are meaningful moves here. Second, the ladder's big steps are
not architectural: engineered features ($+0.0005$), online refit
($+0.0013$), ensembling ($+0.0012$). The architecture rows differ by less
than the operational rows. This pattern --- \textbf{operations beat
architecture} --- recurs in every low-SNR competition we studied.

\begin{warnbox}[: scores that are too good]
On this metric, at this SNR, a validation score above $\sim 0.03$ is not a
breakthrough; it is a leak until proven otherwise. The fastest check:
shift the suspect feature set forward by one timestep and re-score. Genuine
signal degrades gracefully; leaked signal collapses. (The full leakage
audit protocol is in Chapter~\ref{ch:leakage}.) In our campaign this reflex
caught, among others, a rolling statistic computed with a centered window ---
a one-character bug (\code{center=True}) worth a fake $+0.02$.
\end{warnbox}

\section{Metric-driven post-processing}

Because the metric is known in closed form, some post-processing is
mathematically justified rather than superstitious:

\begin{itemize}
\item \textbf{Global rescale.} Fit $\alpha^\star$ on a held-out window
  as above. Values persistently below $1$ indicate overconfident
  magnitudes --- common for models trained with insufficient
  regularization.
\item \textbf{Clipping.} $\hat y \mapsto \mathrm{clip}(\hat y, -c, c)$ with
  $c$ set from the target's realized range (we use $c=5$ on a target whose
  bulk lies within $\pm 3$). The expected cost of clipping true extremes is
  negligible; the protection against pathological rows is not.
\item \textbf{Per-day or per-symbol debiasing --- handle with care.}
  Subtracting each day's running prediction mean sounds attractive
  (the metric's zero-origin punishes bias) but at streaming inference you
  only know the running mean \emph{so far}, and early-day debiasing adds
  variance. We measured no net gain and dropped it.
\end{itemize}

\section{General lessons for any competition}

\begin{enumerate}
\item Read the metric's formula, not its name. ``$R^2$'' here meant
  \emph{zero-mean weighted} $R^2$, and each qualifier changed behavior.
\item Locate the zero point (what does a trivial predictor score?) and the
  ceiling (what did winners of similar contests score?). All progress is
  measured between those posts.
\item Push the metric into training: weights into the loss, weights into
  early stopping, the metric itself as the loss where batch structure
  allows.
\item Derive what post-processing the metric legitimizes (shrinkage,
  clipping, debiasing) and validate each one; do not cargo-cult them.
\item Recompute the metric yourself. Never trust a framework's default
  ``r2'' to be the competition's formula --- ours is eleven characters
  different and those characters are worth points.
\end{enumerate}

\begin{drillbox}
Derive the optimal shrinkage $\alpha^\star$ for the zero-mean weighted
$R^2$: with $\hat y = \alpha f$, expand \eqref{eq:wr2} and minimize over
$\alpha$. Show that using $\alpha=1$ when the true optimum is $\alpha^\star$
costs exactly $\Var_w(f)\,(1-\alpha^\star)^2 / \sum_i w_i y_i^2$ of score.
Then compute $\alpha^\star$ for your current model's validation predictions.
If it is far from $1$, you have found free score.
\end{drillbox}
